\section{Compito Uno: generazione delle griglie}
\label{sec:compito1}

\subsection{Parametri}

Il generatore \code{GeneratoreGriglia.genera(W, H, N, seed)} produce una
griglia di $W \times H$ celle con esattamente $N$ ostacoli in posizioni
casuali. L'utente può configurare:

\begin{itemize}
  \item le dimensioni $W$ e $H$, che devono essere positive;
  \item il numero di ostacoli $N$, con $0 \le N \le W H$; dalla riga di
        comando si impone anche $N < W H$, perché servono almeno una cella
        libera per l'origine e una per la destinazione;
  \item il \emph{seed} del generatore pseudo-casuale.
\end{itemize}

La densità degli ostacoli $N/(WH)$ è il parametro che più influenza il costo
di \cmin{} (\cref{sec:limiti-scala}). Il seed rende ogni griglia
riproducibile: la stessa quaterna $(W,H,N,\mathit{seed})$ produce sempre la
stessa griglia, perché \code{java.util.Random} è un generatore deterministico
a partire dal seed~\cite{random}.

\subsection{Algoritmo}

Le celle vengono estratte uniformemente; una matrice booleana evita di
scegliere due volte la stessa cella. Le celle estratte sono raccolte in una
lista, che viene aggiunta alla griglia con un'unica operazione quando sono
stati scelti esattamente $N$ ostacoli (\cref{alg:genera}).

\begin{algorithm}[H]
\caption{Generazione di una griglia con ostacoli casuali}
\label{alg:genera}
\begin{algorithmic}[1]
\Require dimensioni $W, H > 0$, numero di ostacoli $0 \le N \le WH$, $\mathit{seed}$
\Ensure griglia $W \times H$ con esattamente $N$ ostacoli
\Procedure{GeneraGriglia}{$W, H, N, \mathit{seed}$}
  \State $g \gets$ griglia $W \times H$ con tutte le celle attraversabili
  \State $\mathit{occupata} \gets$ matrice booleana $H \times W$ inizializzata a \textbf{false}
  \State $\mathit{rng} \gets \textsc{Random}(\mathit{seed})$
  \State $L \gets \seq{}$ \Comment{celle scelte come ostacoli}
  \While{$|L| < N$}
    \State $x \gets \mathit{rng}.\textsc{Intero}(W)$; \quad $y \gets \mathit{rng}.\textsc{Intero}(H)$
    \If{\textbf{not} $\mathit{occupata}[y][x]$}
      \State $\mathit{occupata}[y][x] \gets$ \textbf{true}
      \State aggiungi $(x, y)$ a $L$
    \EndIf
  \EndWhile
  \State \Return $g.\textsc{AggiungiOstacoli}(L)$
\EndProcedure
\end{algorithmic}
\end{algorithm}

\paragraph{Costo.} Quando sono già stati collocati $i$ ostacoli, un'estrazione
trova una cella libera con probabilità $(WH - i)/(WH)$. Il numero atteso di
estrazioni è quindi
\[
  \sum_{i=0}^{N-1} \frac{WH}{WH - i},
\]
che resta al più $2N$ finché $N \le WH/2$. L'inizializzazione della
matrice e l'unica vista finale (\cref{sec:immutabilita}) costano $O(WH)$,
quindi il costo complessivo è $O(WH)$ per $N \le WH/2$. Una prima versione
creava una vista per ogni ostacolo, copiando ogni volta la mappa di bit, con
costo $O(N \cdot WH)$: la versione attuale produce le stesse griglie, come
verifica un test dedicato, e rende la generazione trascurabile anche per
griglie molto più grandi di quelle usate nel progetto.

\subsection{Griglie immutabili}
\label{sec:immutabilita}

Aggiungere un ostacolo non modifica la griglia esistente: il metodo
\code{addObstacle} restituisce una nuova \code{GrigliaOstacoliView}, formata
dalla griglia di base e da un \code{BitSet}~\cite{bitset} con tutti gli
ostacoli, quelli della base e quelli aggiunti. Se la griglia di partenza è
già una vista, il suo \code{BitSet} viene copiato ed esteso, così le viste
non formano catene: verificare se una cella è un ostacolo costa $O(1)$ e non
alloca oggetti. Lo stesso \code{BitSet}, mai modificato dopo la
costruzione, è usato come chiave della cache di \cmin{}
(\cref{sec:ottimizzazioni}).

Questa scelta serve soprattutto al Compito Tre. Ogni chiamata ricorsiva di
\cmin{} lavora su una griglia in cui la chiusura della chiamata precedente è
diventata un insieme di ostacoli; i rami ricorsivi fratelli devono però
partire dalla stessa griglia. Con viste immutabili ogni ramo riceve la propria
griglia e nessun ramo vede gli ostacoli aggiunti da un altro.

\subsection{Persistenza in JSON}

\code{LettoreGriglia} salva e carica una griglia come array JSON di righe.
Ogni elemento è lo stato della cella secondo la \cref{tab:statocella}: $0$
per una cella libera, $128$ per un ostacolo.

\begin{lstlisting}
[
  [0, 128,   0],
  [0,   0, 128]
]
\end{lstlisting}

Il caricamento rifiuta JSON non valido, matrici vuote o non rettangolari e
valori fuori dall'intervallo $[0,255]$, con un'eccezione esplicita invece di
un risultato parziale. Nel piano sperimentale ogni griglia generata viene
salvata in \code{risultati/griglie\_json/} e ricaricata prima di eseguire
\cmin{}: i 23 file prodotti sono gli input effettivi delle misure e
permettono di ripetere ogni esperimento sulla stessa griglia.
